\chapter{\texorpdfstring{\texttt{create\_agent}}{create\_agent} and Middleware}
\label{ch:agents}

\begin{quote}\itshape
If you have written an ASP.NET Core middleware pipeline, you have written this
before.
\end{quote}

\noindent
\api{create\_agent} is the shortest path to a working agent, and middleware is
1.0's flagship extensibility mechanism. The analogy to \code{app.Use(async (ctx,
next) => ...)} is exact rather than approximate, which makes this the chapter
where a .NET reader's existing instincts are worth the most. The part worth
slowing down for is where the hooks run: inside the compiled graph, which is why
an agent with a full middleware stack can be dropped into a larger graph as a
single node.

\chapterstub{
\textbf{Argument.} Middleware is not a second runtime bolted onto the agent. The
hooks execute inside the LangGraph graph \api{create\_agent} returns, and that is
what makes composition work --- and what makes Chapter~\ref{ch:stategraph}'s
material apply to agents you did not build by hand.

\textbf{Sections.} (7.1) The full signature, argument by argument: model, tools,
system prompt, middleware, \code{response\_format}, checkpointer, store,
\code{context\_schema}, name. (7.2) The agent loop, stated precisely, and its
termination conditions. (7.3) The six hooks and their ordering:
\code{before\_agent}, \code{before\_model}, \code{wrap\_model\_call},
\code{after\_model}, \code{wrap\_tool\_call}, \code{after\_agent}. (7.4) The
built-in catalogue --- summarisation, human-in-the-loop, PII, model and tool call
limits, model fallback, tool selection --- with a table of when each earns its
place. (7.5) Writing your own: \code{@dynamic\_prompt}, \code{@wrap\_tool\_call},
and the decorator versus class forms. (7.6) Ordering effects, worked: the same
three middleware in two orders, producing two different behaviours. (7.7) The
ASP.NET Core mapping, in full. (7.8) Proving the composition claim by printing
the graph and then embedding the agent as a node.

\textbf{Listings.} A complete agent with five middleware; an audit middleware
written from scratch; a middleware ordering demonstration; the agent-as-node
composition.

\textbf{Diagrams.} \code{agent-loop} --- the model/tool loop as a flowchart with
its exits. \code{middleware-onion} --- the hook order drawn as nested scopes
around one model call. \code{middleware-aspnet} --- the two pipelines side by
side, ASP.NET Core and LangChain.

\textbf{Mini-project.} Stage 04: the agent assembled, with audit, PII redaction,
call limits and an approval gate on the email tool.

\textbf{Source topics covered.} Part~III 3.11--3.13.
}
